\section{实验设置}

\subsection{数据集}

我们在两个自然连续语音聆听数据集上进行实验：法语 MEG 数据集 Pallier2025 和中文 EEG 数据集 ChineseEEG-2，两者的刺激材料均为《小王子》有声书。
训练、验证和测试集按语言篇章划分，所有受试者使用相同划分，测试集篇章在模型训练中完全未见。

ChineseEEG-2\cite{chen2026chineseeeg2} 使用 Passive Listening 条件下 8 名受试者的 128 通道 EEG。
由于实际朗读与原文之间存在停顿、漏读、增读或替换，
我们依据实际播放音频重新核对了词级内容和时间戳。
中文使用 jieba 精确模式分词，关闭 HMM 新词识别；
英文和阿拉伯数字保留原文本组，标点不作为带时间戳的词事件。
第 3、5、25 章用于验证，第 2、9 章用于测试，其余有效章节用于训练。

Pallier2025\cite{openneuro_pallier2025_listening,dascoli2025word} 使用 10 名受试者的 306 个 MEG 传感器（102 个磁强计、204 个平面梯度计），
以 listening run 为划分单位：run-07 用于验证，run-06 用于测试，其余有效 run 用于训练。

有官方预处理结果的数据直接使用该结果，所有信号统一重采样至 $50$~Hz。


\subsection{模型与训练}

本文以 d'Ascoli 等人的词级解码框架为基础\cite{dascoli2025word}。
文本目标由冻结的 BGE-M3 编码器\cite{chen2024bgem3} 得到：每个词单独编码，
取末层首个 token 的 $1024$ 维隐藏状态，并做 $L_2$ 归一化。
训练词库只包含训练材料中的词。

神经端以每个词语音起点后 $1\,\mathrm{s}$ 的信号为输入，
经基于传感器位置的通道映射和残差膨胀卷积网络得到 $\mathbf z_i$，
上下文模型为 Transformer。
训练按第~\ref{sec:revision_training}~节的两阶段进行，
每个数据集使用三个随机种子。
网络结构与超参数见附录表~\ref{tab:model_settings}。

上下文范围实验比较半径 $k\in\{0,1,2,4,8\}$ 的对称窗口，
以及使用片段内全部词的完整上下文模型。
各条件均包含上下文 Transformer，半径 $0$ 即只输入目标词。
可见状态实验固定两个历史词，
将共享滚动模型的三个状态与分别训练的 $[-2,r]$ 模型（$r=0,1,2$）进行比较。
验证集上，对称窗口的平均宏平均 Top-$10$ 在 Pallier2025 中以 $k=2$ 最高，
在 ChineseEEG-2 中以 $k=4$ 最高。
结合局部窗口的验证结果，滚动模型统一取两个历史词窗和最多两个后续词窗，即 $k=R=2$，
以便在两种数据条件下使用相同的回修预算。
检查点和语言权重在验证集上选择，测试前冻结解码配置。

为考察后续词窗的收益是否能由延长目标信号解释，
另从头训练不含跨词上下文、窗口为 $1.28$ 和 $1.50\,\mathrm{s}$ 的单词模型。

\subsection{评价}

固定 $50$ 词候选取训练材料中频率最高的规范化词，按独立材料词事件计频，词频相同时按词典序排序。
词级检索在该冻结集合上，按神经表示与文本表示的余弦相似度排序，
报告 Top-$1$、Top-$10$ 的微平均与宏平均，以宏平均 Top-$10$ 为主指标。
微平均按词位统计，宏平均对真实词类型等权；均匀随机排序时，$50$ 选 $10$ 的机会水平为 $20\%$。
上下文范围比较使用真实词属于候选集合的全部合格查询，
可见状态比较只使用三个状态均可评价的词位。
完整训练词库评价保留词库外目标并计为未命中。

\label{sec:sentence_experiment}
文本序列实验比较表~\ref{tab:strategies} 中的六种策略，语言模型为 Qwen2.5-1.5B 基础版\cite{yang2024qwen25}。
Pallier2025 按原始词标注的 \texttt{sequence\_id} 在各记录内分组，
ChineseEEG-2 根据实际朗读标点和长度限制切分（附录~\ref{sec:supp_settings}）。
包含各受试者的重复呈现后，测试集分别有 $1940$ 和 $808$ 个片段。
语言权重的验证选择及搜索参数见附录~\ref{sec:supp_sentence_settings}。
提示要求输出连续叙述正文，不包含刺激的书名或内容；
法语前缀还原省音形式，中文前缀直接拼接，候选语言评分使用加入候选前后整串文本的对数概率之差。

文本主指标为平均片段 BLEU-1，先在受试者内平均，再对受试者和训练种子等权平均；
WER 和词位准确率按词位汇总，再对训练种子平均。
输入等待按第~\ref{sec:input_wait}~节计算，并按词位等权汇总。
发生文本变化的回修分为错改对、对改错和错改为另一错误三类。
图表中的区间为 $95\%$ 材料块 bootstrap 区间。
